% Part: first-order-logic
% Chapter: syntax-and-semantics
% Section: substitution

\documentclass[../../../include/open-logic-section]{subfiles}

\begin{document}

\olfileid[pt]{fol}{syn}{ext}

\olsection{Extensionalidade}

\begin{explain}
A extensionalidade, às vezes chamada de relevância, pode ser expressa
informalmente da seguinte forma: os únicos fatores que influenciam a
satisfação de !!{formula}~$!A$ em !!a{structure}~$\Struct M$
em relação a uma atribuição de variáveis~$s$ são o tamanho do
!!{domain} e as atribuições feitas por~$\Struct M$ e~$s$ aos
elementos da linguagem que de fato aparecem em~$!A$.

Uma consequência imediata da extensionalidade é que, quando duas
!!{structure}s~$\Struct M$ e~$\Struct M'$ concordam em todos os elementos
da linguagem que aparecem em uma sentença~$!A$ e têm o mesmo
domínio,~$\Struct M$ e~$\Struct M'$ também devem concordar quanto a se
$!A$ em si é verdadeira ou não.
\end{explain}

\begin{prop}[Extensionalidade]
  \ollabel{prop:extensionality}
  Seja $!A$ !!a{formula}, e $\Struct M_1$ e $\Struct M_2$
  !!{structure}s com $\Domain{M_1} = \Domain{M_2}$, e $s$ uma
  atribuição de variáveis sobre $\Domain{M_1} = \Domain{M_2}$. Se
  $\Assign{c}{M_1} = \Assign{c}{M_2}$, $\Assign{R}{M_1}=\Assign{R}{M_2}$,
  e $\Assign{f}{M_1} = \Assign{f}{M_2}$ para todo !!{constant}~$c$,
  símbolo de relação~$R$, e !!{function} $f$ que ocorre em~$!A$, então
  $\Sat{M_1}{!A}[s]$ sse $\Sat{M_2}{!A}[s]$.
\end{prop}

\begin{proof}
  Primeiro prove (por indução sobre~$t$) que, para todo termo,
  $\Value{t}{M_1}[s] = \Value{t}{M_2}[s]$. Então prove a proposição
  por indução sobre~$!A$, fazendo uso da afirmação que acabamos de provar para a
  base da indução (onde $!A$ é atômica).
\end{proof}

\begin{prob}
Faça a prova de \olref[fol][syn][ext]{prop:extensionality} em
detalhe.
\end{prob}

\begin{cor}[Extensionalidade para !!^{sentence}s]
  \ollabel{cor:extensionality-sent}
  Seja $!A$ !!a{sentence} e $\Struct{M_1}$, $\Struct{M_2}$ como em
  \olref{prop:extensionality}. Então $\Sat{M_1}{!A}$ sse $\Sat{M_2}{!A}$.
\end{cor}

\begin{proof}
Segue de \olref{prop:extensionality} por \olref[ass]{cor:sat-sentence}.
\end{proof}

Além disso, o valor de um termo, e se !!a{structure}
satisfaz ou não !!a{formula}, dependem apenas dos valores de seus subtermos.

\begin{prop}\ollabel{prop:ext-terms}
Seja $\Struct M$ !!a{structure}, $t$ e $t'$ termos, e $s$ uma
atribuição de variáveis. Então $\Value{\Subst{t}{t'}{x}}{M}[s] =
\Value{t}{M}[\Subst{s}{\Value{t'}{M}[s]}{x}]$.
\end{prop}

\begin{proof}
Por indução sobre~$t$.
\begin{enumerate}
\item Se $t$ é uma constante, digamos, $t\ident c$, então $\Subst{t}{t'}{x} =
  c$, e $\Value{c}{M}[s] = \Assign{c}{M} =
  \Value{c}{M}[\Subst{s}{\Value{t'}{M}[s]}{x}]$.

\item Se $t$ é uma variável diferente de~$x$, digamos, $t \ident y$, então
  $\Subst{t}{t'}{x} = y$, e $\Value{y}{M}[s] =
  \Value{y}{M}[\Subst{s}{\Value{t'}{M}[s]}{x}]$ já que
  $\varAssign{s}{\Subst{s}{\Value{t'}{M}[s]}{x}}{x}$.

\item Se $t \ident x$, então $\Subst{t}{t'}{x} = t'$. Mas
  $\Value{x}{M}[\Subst{s}{\Value{t'}{M}[s]}{x}] = \Value{t'}{M}[s]$ pela
  definição de~$\Subst{s}{\Value{t'}{M}[s]}{x}$.

\item Se $t \ident \Atom{f}{t_1,\dots,t_n}$ então temos:
\begin{multline*}
  \Value{\Subst{t}{t'}{x}}{M}[s]  = \\
  \begin{aligned}[b]
& = \Value{\Atom{f}{\Subst{t_1}{t'}{x}, \dots, \Subst{t_n}{t'}{x}}}{M}[s]\\
& \qquad    \text{ pela definição de $\Subst{t}{t'}{x}$}\\
& = \Assign{f}{M}(\Value{\Subst{t_1}{t'}{x}}{M}[s], \dots,
    \Value{\Subst{t_n}{t'}{x}}{M}[s])\\
    & \qquad  \text{ pela definição de $\Value{\Atom{f}{\dots}}{M}[s]$}\\
& = \Assign{f}{M}(\Value{t_1}{M}[\Subst{s}{\Value{t'}{M}[s]}{x}], \dots,
   \Value{t_n}{M}[\Subst{s}{\Value{t'}{M}[s]}{x}])\\
& \qquad    \text{ pela hipótese de indução}\\
& = \Value{t}{M}[\Subst{s}{\Value{t'}{M}[s]}{x}]
    \text{ pela definição de $\Value{\Atom{f}{\dots}}{M}[\Subst{s}{\Value{t'}{M}[s]}{x}]$}
  \end{aligned}
\end{multline*}
\end{enumerate}
\end{proof}

\begin{prop}\ollabel{prop:ext-formulas} Seja $\Struct M$
!!a{structure}, $!A$ !!a{formula}, $t'$~um termo, e $s$~uma
atribuição de variáveis. Então $\Sat{M}{\Subst{!A}{t'}{x}}[s]$ sse
$\Sat{M}{!A}[\Subst{s}{\Value{t'}{M}[s]}{x}]$.
\end{prop}

\begin{proof}
Exercício.
\end{proof}

\begin{prob}
Prove \olref[fol][syn][ext]{prop:ext-formulas}
\end{prob}

\begin{explain}
  O ponto de
  \cref{fol:syn:ext:prop:ext-terms,fol:syn:ext:prop:ext-formulas} é
  o seguinte. Suponha que temos um termo $t$ ou !!a{formula}~$!A$ e
  algum termo~$t'$, e queremos saber o valor de $\Subst{t}{t'}{x}$
  ou se $\Subst{!A}{t'}{x}$ é ou não satisfeita em
  !!a{structure}~$\Struct M$ em relação a uma atribuição de variáveis~$s$.
  Então podemos ou realizar a substituição primeiro e depois considerar
  o valor ou a satisfação em relação a $\Struct{M}$ e~$s$, ou
  podemos primeiro determinar o valor~$m = \Value{t'}{M}[s]$ de $t'$ em
  $\Struct{M}$ em relação a~$s$, mudar a atribuição de variáveis
  para~$\Subst{s}{m}{x}$ e então considerar o valor de~$t$ em
  $\Struct{M}$ e~$\Subst{s}{m}{x}$, ou se
  $\Sat{M}{!A}[\Subst{s}{m}{x}]$.
  \Cref{fol:syn:ext:prop:ext-terms,fol:syn:ext:prop:ext-formulas}
  garantem que a resposta será a mesma, independentemente de como o façamos.
\end{explain}

\end{document}
